\clearpage
\fancyhead[L]{\small HELLER FLECK IM MOND.}
\SourceLeaf{441}
\begin{center}
{\large [HELLER FLECK IM MOND.]}
\end{center}

\begin{center}
\noindent\rule{\linewidth}{0.4pt}

Göttingische gelehrte Anzeigen. Nr. 46. S. 449..452. 22. März 1821.

\noindent\rule{\linewidth}{0.4pt}
\end{center}

Über eine Erscheinung, die Dr. \textsc{Olbers} am 5. Februar dieses Jahres auf dem dunklen Teil der Mondoberfläche beobachtet hat, teilen wir hier einen Auszug aus seinem Brief vom 27. Februar an Hofrat \textsc{Gauss} mit. Das tun wir umso lieber, als die überaus einfallsreichen Vermutungen, die dieser geniale Astronom über das an sich schon ebenso seltene wie bemerkenswerte Phänomen anfügt, die größte Aufmerksamkeit verdienen und zu einer ganz neuen und fruchtbaren Sichtweise führen.

Am 5. Februar (schreibt Dr. \textsc{Olbers}) habe ich die Erscheinung auf dem Mond gesehen, die man einen Mondvulkan genannt hat. Was ich darüber sogleich in mein Tagebuch geschrieben habe, lautet wörtlich wie folgt. `Am Fünften war es sehr klar, aber es gab schon Mondschein. Auf dem dunklen Teil des Mondes habe ich das Phänomen, das man für einen brennenden Vulkan auf dem Mond gehalten hat, noch nie so deutlich und auffällig gesehen wie an diesem Abend. Wie gewöhnlich schien es sich im Aristarch zu befinden. Es war klein, aber ganz \emph{auffallend} heller als der übrige, von der Sonne nicht beleuchtete Teil des Mondes. Es sah ganz wie ein Stern aus, und zwar genau wie ein Fixstern dritter Größe nordöstlich vom Mond.'

Da es am 6. Februar bedeckt war, hat Dr. \textsc{Olbers} nicht weiter nach dieser Erscheinung Ausschau gehalten. Bald darauf meldeten jedoch englische Zeitungen, Kapitän \textsc{Kater} habe der Royal Society in London am 7. Februar von einem Vulkan berichtet, den er auf dem Mond gesehen habe. Durch fortgesetzte Beobachtungen habe er sich tatsächlich davon überzeugt, dass der Vulkan gerade ausbreche.

Es scheint also (fährt Dr. \textsc{Olbers} fort), dass Herr \textsc{Kater} dieselbe Erscheinung gesehen hat, die auch mir am 5. Februar auffiel, sie aber weiter verfolgt hat. Ich weiß zwar nicht, was ihn davon überzeugt hat, dass es wirklich ein brennender Vulkan war. Doch nach allem, was wir über die Beschaffenheit des Mondes und seine so zweifelhafte Atmosphäre wissen, scheint ein brennender Vulkan fast unmöglich. `Vielmehr glaube ich, dass sich die Erde tatsächlich in der ebenen, merklich glatten, beinahe wie poliert wirkenden Seitenwand einer großen, zum Aristarch gehörenden Felsklippe gespiegelt hat.' Das so gespiegelte Bild eines Teils der Erde musste sehr viel heller sein als alles Übrige, was allein von der Erde beleuchtet wurde. Denn dieses streut das Erdlicht in alle Richtungen, während jenes es nur in \emph{eine} Richtung zurückwirft. Selbst wenn jene unvollkommene Spiegelung nur ein Zehntel des Erdlichts zurückwarf (da unsere wirklichen Spiegel etwa die Hälfte des auf sie fallenden Lichts zurückwerfen) und die Seitenwand nur $2''$ im Durchmesser hatte, konnte sie doch so hell wie ein Stern sechster Größe erscheinen. --- Mit dieser Vorstellung lässt sich \emph{erstens} erklären, warum wir die vulkanartigen Erscheinungen immer nur an bestimmten Stellen des Mondes sehen. \emph{Zweitens} lässt sich erklären, warum wir sie nur selten und nicht in jeder Lunation zu Gesicht bekommen. Die Libration muss nämlich bis auf etwa $2^\circ$ dieselbe sein. --- Dass es solche Seitenwände von Mondklippen geben kann, die das Licht mehr oder weniger wie ein Spiegel zurückwerfen, lässt sich wohl kaum bezweifeln. Auf der Erde mag es große Gletscherflächen geben, die ebenfalls als unvollkommene Spiegel Licht zurückwerfen können. Ich führe das nur als Analogie an, denn Gletscher auf dem Mond sind ebenso
\SourceLeaf{442}%
unwahrscheinlich wie brennende Vulkane. Doch es scheint durchaus vorstellbar, dass unter den Gebirgen, die sich wahrscheinlich auch auf dem Mond nach Kristallisationsgesetzen gebildet haben, einzelne ebene, glatte, beinahe wie poliert wirkende Seitenflächen haben. Vielleicht gibt es davon viele auf dem Mond, aber sie dürften selten gerade so liegen, dass sie uns bei einer bestimmten Libration genau das Bild der Erde zurückspiegeln können. Auch die Sonne scheint sich zuweilen in ähnlichen Klippenwänden auf dem Mond zu spiegeln. Noch vor etwa 8 Wochen sah ich im Mare Imbrium jenseits der Lichtgrenze zwei von der Sonne beschienene Berggipfel. Sie waren so ungewöhnlich hell, funkelnd und dem Sirius ähnlich, dass ich es für unmöglich hielt, hier bloß nach den gewöhnlichen Streuungsgesetzen zurückgeworfenes Sonnenlicht anzunehmen. Ich kann nicht mit Sicherheit sagen, ob einer dieser Berge vielleicht der ebenfalls im Mare Imbrium gelegene Lahire (nach \textsc{Schröter}) war, an dem \textsc{Schröter} ganz ähnliche Erscheinungen beobachtet hat.

In einem späteren Brief bemerkt Dr. \textsc{Olbers} noch, dass sich diese Hypothese über die Ursache des Phänomens leicht prüfen lassen werde. Falls sie die wahre Erklärung liefere, müsse bei derselben Libration immer wieder dieselbe Erscheinung auftreten. Felswände, die ein Bild der Erde oder der Sonne mehr oder weniger unvollkommen zurückspiegeln können, seien auf dem Mond umso eher vorstellbar, da dort die äußere Oberfläche der Gebirge und Klippen wahrscheinlich nicht wie auf der Erde durch atmosphärische Einwirkung verwittere. Die spiegelnde Klippenwand brauche auch nicht ganz eben zu sein, sofern nur die spiegelnden Teile in parallelen Ebenen liegen, wie es bei Bergen, die sich nach Kristallisationsgesetzen gebildet haben, leicht der Fall sein könne. Dr. \textsc{Olbers} erinnert hier an unsere Basaltberge. Deren einzelne große Kristalle könnten zusammen für ein entferntes Auge noch sehr wohl ein unvollkommenes Sonnenbild zurückspiegeln, wenn ihre Oberflächen nicht längst durch Luft, Dünste, Regen usw. ihre wahrscheinlich ursprünglich vorhandene Politur und Glätte verloren hätten.

\Needspace{7\baselineskip}Am 6. März, als die Nachtseite des Mondes ausgezeichnet zu sehen war, konnte Dr. \textsc{Olbers} mit seinem Fernrohr von \textsc{Dollond} alle Flecken sehr deutlich erkennen, etwa Grimaldi, Copernicus, Kepler, Maeilius, Menelaus usw. Aristarch hob sich wieder von allen anderen ab, anscheinend sogar \emph{mehr als gewöhnlich}. Doch so \emph{hell} und so \emph{fixsternähnlich} wie am 5. Februar fand ihn Dr. \textsc{Olbers} diesmal nicht.
